\documentclass[11pt,a4paper]{article}

% ============================================================
%  Inventario_Proyecto.tex
%  UrbanFlow CDMX - Inventario Completo del Proyecto
%  Carlos Armando López Encino
%  Diplomado en Ciencia de Datos - FES Acatlán - UNAM
%  Abril 2026
% ============================================================

% --- Codificación e idioma ---
\usepackage[utf8]{inputenc}
\usepackage[T1]{fontenc}
\usepackage[provide=*,main=spanish]{babel}
\usepackage{lmodern}
\usepackage{microtype}

% --- Geometría ---
\usepackage[a4paper,margin=2.5cm,headheight=14pt]{geometry}

% --- Colores y cajas ---
\usepackage{xcolor}
\definecolor{navy}{RGB}{15,36,72}
\definecolor{amber}{RGB}{212,149,35}
\definecolor{lightgray}{RGB}{240,240,240}
\definecolor{rulegray}{RGB}{180,180,180}

% --- Tipografía de encabezados ---
\usepackage{titlesec}
\titleformat{\section}
  {\normalfont\Large\bfseries\color{navy}}
  {\thesection}{1em}{}
\titleformat{\subsection}
  {\normalfont\large\bfseries\color{navy!85!black}}
  {\thesubsection}{1em}{}

% --- Tablas ---
\usepackage{array}
\usepackage{booktabs}
\usepackage{tabularx}
\usepackage{longtable}
\newcolumntype{L}[1]{>{\raggedright\arraybackslash}p{#1}}

% --- Listas ---
\usepackage{enumitem}
\setlist[itemize]{leftmargin=1.4em,itemsep=2pt,topsep=3pt}

% --- Cabeceras y pies ---
\usepackage{fancyhdr}
\pagestyle{fancy}
\fancyhf{}
\renewcommand{\headrulewidth}{0.4pt}
\renewcommand{\footrulewidth}{0pt}
\fancyhead[L]{\small\color{navy}UrbanFlow CDMX \(|\) Inventario}
\fancyhead[R]{\small\color{navy}Diplomado en Ciencia de Datos}
\fancyfoot[C]{\thepage}

% --- Hipervínculos ---
\usepackage[hidelinks,colorlinks=true,linkcolor=navy,urlcolor=navy,citecolor=navy]{hyperref}

% --- Símbolos matemáticos ---
\usepackage{amsmath,amssymb}

% --- Código fuente ---
\usepackage{listings}
\lstset{
  basicstyle=\ttfamily\small,
  frame=single,
  rulecolor=\color{rulegray},
  backgroundcolor=\color{lightgray},
  breaklines=true,
  columns=fullflexible,
  xleftmargin=0.5em,
  xrightmargin=0.5em
}

% --- Espaciado ---
\setlength{\parskip}{0.5em}
\setlength{\parindent}{0pt}

% ============================================================
\begin{document}

% ---------- Portada ----------
\begin{titlepage}
\centering
\vspace*{3cm}

{\color{navy}\rule{\textwidth}{1.2pt}}\\[0.8em]
{\Huge\bfseries\color{navy} UrbanFlow CDMX}\\[0.5em]
{\Large\color{navy} Inventario Completo del Proyecto}\\[0.3em]
{\large\itshape\color{amber} Qué se hizo \(\cdot\) Con qué herramientas \(\cdot\) En qué orden}\\[0.3em]
{\color{navy}\rule{\textwidth}{1.2pt}}\\[3cm]

{\Large\bfseries Carlos Armando López Encino}\\[0.4em]
{\large Diplomado en Ciencia de Datos}\\[0.2em]
{\large FES Acatlán \(\cdot\) UNAM}\\[1.5cm]

{\large Proyecto Integrador Final}\\[0.3em]
{\large Abril 2026}

\vfill
\end{titlepage}

% ---------- Propósito ----------
\section*{Propósito de este documento}
\addcontentsline{toc}{section}{Propósito de este documento}

Este documento es un inventario exhaustivo de todo lo que se construyó durante el proyecto integrador del Diplomado en Ciencia de Datos: herramientas, bibliotecas, fuentes de datos, decisiones arquitectónicas, entregables y recursos consultados. Está organizado cronológicamente por fase de trabajo para que refleje el orden real en que fueron apareciendo las piezas del proyecto.

No es un documento de marketing ni un pitch; es un registro técnico interno del esfuerzo realizado, útil para cualquier persona que quiera reproducir el trabajo, continuar el desarrollo, o simplemente entender el alcance del proyecto.

\textit{Nota: las fechas exactas de cada fase están marcadas con \texttt{[POR CONFIRMAR]} en los lugares donde el autor debe completarlas antes de entregar este documento al profesor.}

% ============================================================
\section{Fase 0 --- Concepción y planteamiento \texttt{[POR CONFIRMAR: nov--dic 2025]}}

\subsection{Definición del problema}
\begin{itemize}
  \item Identificación del tráfico en la ZMVM como problema no resuelto por sistemas de navegación convencionales
  \item Revisión del TomTom Traffic Index 2024 y confirmación de que CDMX es la ciudad \#1 en congestión mundial
  \item Definición del gap: Google Maps/Waze entregan estimaciones puntuales; ningún sistema entrega bandas de incertidumbre específicas del contexto mexicano
\end{itemize}

\subsection{Elección del enfoque metodológico}
\begin{itemize}
  \item Decisión de usar cadenas de Markov de tiempo discreto como modelo de dinámica de tráfico
  \item Decisión de complementar con simulación Monte Carlo para cuantificar incertidumbre
  \item Justificación académica: combina dos temas del diplomado (procesos estocásticos + métodos de simulación) en un problema realista
\end{itemize}

\subsection{Fuentes consultadas en la fase inicial}
\begin{itemize}
  \item TomTom Traffic Index 2024 --- \url{https://www.tomtom.com/traffic-index/}
  \item INEGI --- Encuesta Origen-Destino en Hogares ZMVM 2017
  \item DataMéxico (Secretaría de Economía) --- perfil económico de Valle de México
  \item Mordor Intelligence --- Mexico Last Mile Delivery Market Report 2025--2030
  \item Norris, J. R. (1997). \textit{Markov Chains}. Cambridge University Press
  \item Rubinstein \& Kroese (2016). \textit{Simulation and the Monte Carlo Method}. Wiley
\end{itemize}

% ============================================================
\section{Fase 1 --- Arquitectura del sistema \texttt{[POR CONFIRMAR: dic 2025]}}

\subsection{Decisiones arquitectónicas principales}
\begin{itemize}
  \item Separación clara entre motor estocástico (Python puro vectorizado) y capa conversacional (agente LLM)
  \item Estructura de directorios estilo \texttt{src/}:
\end{itemize}

\begin{lstlisting}
UrbanFlow_CDMX/
|-- src/
|   |-- simulation/    (markov_chain.py, monte_carlo.py)
|   |-- data_sources/  (tomtom_client.py, weather_client.py, c5_client.py)
|   |-- agent/         (agent.py, perturbaciones.py, tools.py)
|   +-- core/          (schemas.py Pydantic)
|-- app/
|   +-- streamlit_app.py
|-- tests/             (pytest)
|-- notebooks/         (EDA)
|-- data/              (caches parquet)
+-- .env
\end{lstlisting}

\begin{itemize}
  \item Elección de Anthropic API sobre OpenAI por mejor razonamiento sobre tablas probabilísticas
  \item Elección de Streamlit sobre Gradio por integración más madura con Folium (mapas interactivos)
\end{itemize}

\subsection{Stack tecnológico inicial}
\begin{center}
\begin{tabularx}{\textwidth}{L{3.2cm} L{3.2cm} X}
\toprule
\textbf{Capa} & \textbf{Herramienta} & \textbf{Rol} \\
\midrule
Cómputo numérico      & NumPy 1.26      & Álgebra lineal vectorizada (Markov, Monte Carlo) \\
Datos tabulares       & Pandas 2.2      & DataFrames de tráfico, clima, rutas de validación \\
Geoespacial           & GeoPandas       & Proyección EPSG:6372 para cálculos de distancia \\
Mapas                 & Folium          & Renderizado interactivo de rutas sobre tiles CDMX \\
Validación esquemas   & Pydantic 2.x    & Structured outputs y contratos tipados \\
Frontend              & Streamlit 1.56  & Dashboard interactivo con mapa + chat \\
Agente LLM            & anthropic SDK   & Cliente Python para Claude API \\
HTTP asíncrono        & httpx           & Llamadas concurrentes a APIs externas \\
Testing               & pytest          & 443 tests automatizados \\
Gestión deps          & pip-tools       & Lockfile reproducible en requirements.txt \\
Entorno               & python-dotenv   & Carga de API keys desde .env \\
\bottomrule
\end{tabularx}
\end{center}

% ============================================================
\section{Fase 2 --- Fuentes de datos e ingesta \texttt{[POR CONFIRMAR: ene 2026]}}

\subsection{Cinco fuentes de datos integradas}
\begin{center}
\begin{tabularx}{\textwidth}{L{3.2cm} L{3cm} X}
\toprule
\textbf{Fuente} & \textbf{Tipo} & \textbf{Uso en el sistema} \\
\midrule
TomTom Traffic Stats API & REST, tiempo real & Velocidad actual y de flujo libre por coordenada. Inferencia del estado inicial $s_0$ de la cadena de Markov. Cuota gratuita: 2{,}500 req/día. \\
TomTom Routing API       & REST, por consulta & Distancia real por carretera y geometría de ruta (polilínea). Ground truth para validación empírica. \\
OpenWeatherMap API       & REST, cada 10 min  & Temperatura, precipitación, condición climática. Factor $\phi$ sobre velocidades. Cuota gratuita: 1{,}000 req/día. \\
C5 CDMX Histórico        & CSV desde portal CKAN & Serie histórica de incidentes viales 2018--2024 ($\sim$1M registros). Calibración de la matriz de transición $\hat{P}$ y densidad horaria. \\
Google Maps Distance Matrix & REST, por consulta & Validación cruzada del ETA, geocodificación de direcciones del usuario. \\
\bottomrule
\end{tabularx}
\end{center}

\subsection{Portal de datos abiertos CDMX}
\begin{itemize}
  \item \url{https://datos.cdmx.gob.mx/dataset/incidentes-viales-c5} --- dataset C5 principal
  \item CKAN API: \url{https://datos.cdmx.gob.mx/api/3/action/package_show}
\end{itemize}

\subsection{Módulos de cliente implementados}
\begin{itemize}
  \item \texttt{src/data\_sources/tomtom\_client.py} --- wrapper de las dos APIs de TomTom
  \item \texttt{src/data\_sources/weather\_client.py} --- cliente OpenWeatherMap con mapeo de condiciones climáticas a factor $\phi$
  \item \texttt{src/data\_sources/c5\_client.py} --- descarga y limpieza del CSV histórico de C5 CDMX
  \item \texttt{src/data\_sources/gmaps\_client.py} --- geocodificación y Distance Matrix
\end{itemize}

\subsection{Retos encontrados y cómo se resolvieron}
\begin{itemize}
  \item \textbf{Cuotas API limitadas:} Se implementó un sistema de caché local con \texttt{requests-cache} para evitar llamadas repetidas durante desarrollo
  \item \textbf{Indisponibilidad ocasional de TomTom Routing:} Se implementó el fallback empírico $d_{\text{carretera}} \approx d_{\text{haversine}} \times 1.4$, calibrado sobre 500 pares O-D aleatorios
  \item \textbf{Nombres de columnas distintos por año en C5:} Se creó un mapeo \texttt{MAPEO\_COLUMNAS} que normaliza nombres a un esquema único
\end{itemize}

% ============================================================
\section{Fase 3 --- Análisis exploratorio de datos \texttt{[POR CONFIRMAR: ene--feb 2026]}}

\subsection{Notebook EDA v3}
Archivo: \texttt{notebooks/EDA\_UrbanFlow\_CDMX\_v3.ipynb}
\begin{itemize}
  \item Limpieza y tipado de $\sim$1 millón de registros del C5 CDMX
  \item Análisis de valores faltantes con umbrales de 5\,\% y 10\,\%
  \item Distribución geográfica (hexbin map) de incidentes sobre el bounding box ZMVM
  \item Distribución temporal: hora del día, día de la semana, mes, heatmap hora$\times$día
  \item Identificación de los dos picos diarios: matutino 7--9\,h y vespertino 17--20\,h
  \item Top alcaldías por incidentes: Iztapalapa, Cuauhtémoc, Gustavo A. Madero
  \item Proxy de congestión por densidad horaria normalizada a $[0.35, 1.0]$
\end{itemize}

\subsection{Bibliotecas adicionales de visualización}
\begin{itemize}
  \item \texttt{matplotlib} --- gráficos base
  \item \texttt{seaborn} --- heatmaps y estilos
  \item \texttt{scipy.stats} --- distribuciones y tests
\end{itemize}

% ============================================================
\section{Fase 4 --- Motor estocástico (Markov + Monte Carlo) \texttt{[POR CONFIRMAR: feb 2026]}}

\subsection{Clase \texttt{MarkovTrafficChain}}
Archivo: \texttt{src/simulation/markov\_chain.py}
\begin{itemize}
  \item Interfaz estilo scikit-learn: \texttt{fit()} $\rightarrow$ \texttt{predict\_distribution()} $\rightarrow$ \texttt{simulate()}
  \item Estimación MLE con suavizado de Laplace ($\varepsilon = 10^{-6}$) para evitar probabilidades cero
  \item Cálculo del estado estacionario $\pi^{*}$ vía descomposición espectral de $P^{\top}$
  \item Soporte para 3 estados: Fluido / Lento / Congestionado
  \item Matriz calibrada para el corredor Insurgentes Sur en pico AM
\end{itemize}

\subsection{Clase \texttt{MonteCarloEngine}}
Archivo: \texttt{src/simulation/monte\_carlo.py}
\begin{itemize}
  \item Simulación vectorizada sobre $N$ trayectorias (10{,}000 por defecto)
  \item Método de la transformada inversa vectorizada para muestrear estados de Markov en paralelo
  \item Muestreo de velocidades por estado usando Normal truncada
  \item Interpolación lineal sub-paso para precisión de tiempo de llegada
  \item Cálculo de percentiles $P_{10}$, $P_{50}$, $P_{90}$ y métrica derivada IC
  \item Latencia: $<$500\,ms por consulta en CPU sin GPU
\end{itemize}

\subsection{Retos técnicos}
\begin{itemize}
  \item \textbf{Vectorización:} la versión naive con bucles Python tomaba $\sim$30 segundos por consulta. La vectorización por transformada inversa la redujo a $<$500\,ms (60$\times$ mejora)
  \item \textbf{Estabilidad numérica:} se agregó normalización después de cada operación sobre la matriz para evitar derivas acumuladas
\end{itemize}

\subsection{Cobertura de tests}
\begin{itemize}
  \item \texttt{tests/test\_markov\_chain.py} --- propiedades de estocasticidad, convergencia estacionaria, predicción en $k$ pasos
  \item \texttt{tests/test\_monte\_carlo.py} --- convergencia del estimador, reproducibilidad con seed, rangos de bandas
  \item \texttt{tests/test\_pipeline.py} --- integración end-to-end desde ratio de congestión hasta bandas $P_{10}/P_{50}/P_{90}$
  \item \textbf{Total: 443 tests pasando al 100\,\%}
\end{itemize}

% ============================================================
\section{Fase 5 --- Catálogo de perturbaciones contextuales \texttt{[POR CONFIRMAR: feb--mar 2026]}}

\subsection{Motivación}
TomTom y OpenWeatherMap no pueden capturar eventos como cierres del Metro, manifestaciones o eventos masivos hasta que ya pasaron. Se construyó un catálogo de perturbaciones que modifica la matriz de Markov \textit{antes} de ejecutar la simulación.

\subsection{Formulación matemática}
Función de perturbación por factor multiplicativo:
\[
p'_{i,2} = \min(f \cdot p_{i,2},\ p_{\max}), \qquad p_{\max} = 0.92
\]
con renormalización automática de la fila para preservar estocasticidad.

\subsection{Catálogo implementado}
Archivo: \texttt{src/agent/perturbaciones.py}
\begin{itemize}
  \item \textbf{Infraestructura:} Cierre Línea 1 Metro ($f = 1.25$), Cierre Línea 12 Metro ($f = 1.20$)
  \item \textbf{Eventos masivos:} Concierto en Centro Banamex ($f = 1.15$), Partido en Estadio Azteca ($f = 1.22$), Concierto en Foro Sol ($f = 1.20$)
  \item \textbf{Manifestaciones:} Marcha 9 de marzo ($f = 1.30$), Marcha magisterial CNTE ($f = 1.27$)
  \item \textbf{Festivos alto impacto:} 15 de septiembre --- Grito ($f = 1.35$), 16 de septiembre --- Desfile ($f = 1.30$)
  \item \textbf{Festivos bajo impacto:} Navidad / Año Nuevo ($f = 0.55$)
\end{itemize}
\textbf{Total: 11 eventos} calibrados empíricamente contra datos históricos del C5 CDMX.

% ============================================================
\section{Fase 6 --- Agente conversacional VialAI \texttt{[POR CONFIRMAR: mar 2026]}}

\subsection{Arquitectura del agente}
\begin{itemize}
  \item Patrón \textit{tool use loop} estándar de la API de Anthropic
  \item Modelo: \texttt{claude-sonnet-4-5} (o versión disponible en el momento)
  \item Integración con Streamlit via \texttt{st.session\_state} para mantener historial
\end{itemize}

\subsection{Structured Outputs con Pydantic}
Archivo: \texttt{src/core/schemas.py}
\begin{itemize}
  \item \texttt{RespuestaTomTom} --- valida respuesta de TomTom Traffic Stats
  \item \texttt{RespuestaClima} --- valida respuesta de OpenWeatherMap
  \item \texttt{PrediccionViaje} --- resultado estructurado del motor ($P_{10}/P_{50}/P_{90}$, IC, semáforo)
  \item \texttt{PerturbacionActiva} --- evento contextual con tipo, radio, duración
\end{itemize}

\subsection{Tres herramientas expuestas al agente}
\begin{itemize}
  \item \texttt{predecir\_tiempo\_viaje(origen, destino, hora)} $\rightarrow$ \texttt{PrediccionViaje}
  \item \texttt{consultar\_trafico\_ahora(zona)} $\rightarrow$ Estado actual + ratio de congestión
  \item \texttt{verificar\_perturbaciones(zona, hora)} $\rightarrow$ Lista de \texttt{PerturbacionActiva}
\end{itemize}

\subsection{Reto: Bug de carga de \texttt{.env}}
\begin{itemize}
  \item \textbf{Síntoma:} \texttt{ANTHROPIC\_API\_KEY} no cargaba cuando Streamlit instanciaba el agente
  \item \textbf{Causa:} \texttt{load\_dotenv()} sin argumentos usa \texttt{cwd} como referencia, pero Streamlit corre desde un directorio distinto al del módulo
  \item \textbf{Solución:} ruta absoluta derivada del propio archivo vía \texttt{Path(\_\_file\_\_).resolve().parents[2]}
\end{itemize}

% ============================================================
\section{Fase 7 --- Interfaz Streamlit \texttt{[POR CONFIRMAR: mar 2026]}}

\subsection{Componentes de la UI}
\begin{itemize}
  \item Mapa interactivo Folium con selección de origen/destino por click
  \item Widget de clima en tiempo real desde OpenWeatherMap
  \item Gauge del Índice de Confiabilidad con semáforo verde/amarillo/rojo
  \item Histograma de la distribución Monte Carlo
  \item Sidebar con chat de VialAI
  \item Toggles para activar perturbaciones del catálogo
\end{itemize}

\subsection{Archivo principal}
\texttt{app/streamlit\_app.py}

\subsection{Dependencias adicionales}
\begin{itemize}
  \item \texttt{streamlit-folium} --- integración mapa-Streamlit
  \item \texttt{plotly} --- gauges interactivos del IC
\end{itemize}

% ============================================================
\section{Fase 8 --- Validación empírica \texttt{[POR CONFIRMAR: mar 2026]}}

\subsection{Dataset de validación}
30 rutas representativas muestreadas entre 9 polos de actividad de la ZMVM: Polanco, Santa Fe, Reforma, Roma-Condesa, Insurgentes Sur, Coyoacán, Aeropuerto, Centro Histórico, Naucalpan.

\subsection{Métricas reportadas}
\begin{itemize}
  \item \textbf{MAPE del $P_{50}$ (VialAI): 11.9\,\%} --- bajo el umbral objetivo del 15\,\%
  \item MAPE del ETA (TomTom Routing): 19.0\,\% --- baseline competitivo
  \item \textbf{VialAI supera a TomTom por 7.1 puntos porcentuales}
  \item Cobertura empírica de la banda $P_{10}$--$P_{90}$: 93.3\,\% (nominal 80\,\%)
  \item MAE del $P_{50}$: 4.8 minutos
\end{itemize}

\subsection{Casos testigo}
\begin{itemize}
  \item \textbf{Caso ganado:} Polanco $\rightarrow$ Aeropuerto (08:15). TomTom ETA: 38\,min. VialAI $P_{50}$: 40\,min. Real: 47\,min (dentro de $P_{90}$). TomTom habría subestimado en 9 minutos.
  \item \textbf{Caso perdido:} Coyoacán $\rightarrow$ Santa Fe (09:30). Real: 95\,min debido a manifestación no anunciada. Fuera de la banda $P_{90}$ de VialAI. Este caso motivó la construcción del catálogo de perturbaciones.
\end{itemize}

% ============================================================
\section{Fase 9 --- Documentación final \texttt{[abril 2026]}}

\subsection{Artículo técnico}
\begin{itemize}
  \item Archivo: \texttt{UrbanFlow\_CDMX\_Articulo\_v2.pdf} (23 páginas)
  \item Fuente: \LaTeX{} con \texttt{tcolorbox}, \texttt{TikZ}, \texttt{listings}, \texttt{microtype}
  \item 12 secciones: motivación, arquitectura, datos, Markov, Monte Carlo, perturbaciones, agente, validación, estrategia B2B, conclusiones
  \item 4 \textit{challengebox} documentando retos técnicos reales
  \item 4 \textit{insightbox} con hallazgos vendibles
  \item 17 referencias bibliográficas reales
\end{itemize}

\subsection{Notebook Colab reproducible}
\begin{itemize}
  \item Archivo: \texttt{UrbanFlow\_CDMX\_Colab.ipynb} (38 celdas)
  \item Datos cacheados inline: matriz de transición + 30 rutas de validación
  \item Corre en Google Colab sin credenciales ni dependencias externas
  \item Reproduce los resultados de validación reportados en el artículo
\end{itemize}

\subsection{Presentación ejecutiva}
\begin{itemize}
  \item Archivo: \texttt{UrbanFlow\_CDMX\_Presentacion.pptx} (16 slides)
  \item Generada con \texttt{pptxgenjs} (Node.js) + iconos de \texttt{react-icons}
  \item Paleta \textit{Midnight Executive} (navy + ámbar)
  \item Estructura narrativa: portada $\rightarrow$ problema $\rightarrow$ metodología $\rightarrow$ demo $\rightarrow$ validación $\rightarrow$ estrategia $\rightarrow$ cierre
\end{itemize}

\subsection{Documentos complementarios}
\begin{itemize}
  \item \texttt{VialAI\_Pitch\_B2B.pdf} --- one-pager ejecutivo de 4 páginas para logística B2B
  \item \texttt{guion\_video.md} --- guion del video de 5 minutos con timestamps
  \item \texttt{checklist\_demo.md} --- checklist operativo para la demo en vivo
\end{itemize}

% ============================================================
\section{Herramientas, bibliotecas y recursos utilizados}

\subsection{Lenguajes de programación}
\begin{itemize}
  \item \textbf{Python 3.14} --- lenguaje principal del motor, agente y pipeline
  \item \textbf{JavaScript (Node.js 22)} --- generación de la presentación con \texttt{pptxgenjs}
  \item \textbf{\LaTeX} --- artículo, pitch B2B e inventario
  \item \textbf{Markdown} --- guion de video, checklist de demo, README
  \item \textbf{Bash / PowerShell} --- scripts de automatización
\end{itemize}

\subsection{Bibliotecas Python (\texttt{requirements.txt})}
\begin{center}
\begin{tabularx}{\textwidth}{L{3.2cm} L{2cm} X}
\toprule
\textbf{Biblioteca} & \textbf{Versión} & \textbf{Uso} \\
\midrule
numpy            & 1.26+ & Álgebra lineal vectorizada \\
pandas           & 2.2+  & DataFrames \\
scipy            & 1.11+ & Distribuciones, tests estadísticos \\
scikit-learn     & 1.3+  & Utilidades de validación \\
matplotlib       & 3.8+  & Visualizaciones base \\
seaborn          & 0.13+ & Heatmaps y estilos \\
geopandas        & 0.14+ & Proyecciones geoespaciales \\
folium           & 0.15+ & Mapas interactivos \\
streamlit        & 1.56+ & Dashboard \\
streamlit-folium & 0.18+ & Integración mapa-dashboard \\
plotly           & 5.18+ & Gauges interactivos \\
pydantic         & 2.5+  & Validación de schemas \\
anthropic        & 0.25+ & Cliente API Claude \\
requests         & 2.31+ & HTTP sincrónico \\
httpx            & 0.26+ & HTTP asíncrono \\
python-dotenv    & 1.0+  & Carga de variables de entorno \\
pytest           & 8.0+  & Testing \\
pytest-cov       & 4.1+  & Cobertura de tests \\
requests-cache   & 1.1+  & Caché local de respuestas API \\
\texttt{[POR CONFIRMAR]} & --- & Otras dependencias del \texttt{requirements.txt} real \\
\bottomrule
\end{tabularx}
\end{center}

\subsection{APIs externas consumidas}
\begin{itemize}
  \item TomTom Traffic Stats API --- \url{https://developer.tomtom.com/traffic-api}
  \item TomTom Routing API --- \url{https://developer.tomtom.com/routing-api}
  \item OpenWeatherMap API --- \url{https://openweathermap.org/api}
  \item Google Maps Distance Matrix API --- \url{https://developers.google.com/maps/documentation/distance-matrix}
  \item Anthropic Messages API --- \url{https://docs.anthropic.com/en/api/messages}
  \item CKAN API (datos.cdmx.gob.mx) --- \url{https://datos.cdmx.gob.mx/api/3}
\end{itemize}

\subsection{Fuentes de datos públicos}
\begin{itemize}
  \item C5 CDMX Incidentes Viales --- \url{https://datos.cdmx.gob.mx/dataset/incidentes-viales-c5}
  \item INEGI Encuesta Origen-Destino ZMVM 2017 --- \url{https://www.inegi.org.mx}
  \item DataMéxico (Sec. Economía) --- \url{https://www.economia.gob.mx/datamexico/es/profile/geo/valle-de-mexico}
  \item SEMOVI aforos vehiculares --- Gobierno CDMX, datasets abiertos
  \item TomTom Traffic Index 2024 --- \url{https://www.tomtom.com/traffic-index/mexico-city-traffic/}
\end{itemize}

\subsection{Software de desarrollo}
\begin{itemize}
  \item \textbf{Visual Studio Code} con extensiones: Python, Pylance, Jupyter, LaTeX Workshop
  \item \textbf{Claude Code} (Anthropic) --- asistencia de programación para refactors y debugging
  \item \textbf{Git + GitHub} --- control de versiones
  \item \textbf{Python venv} --- entorno virtual aislado
  \item \textbf{TeX Live} --- distribución \LaTeX
  \item \textbf{Node.js 22} --- runtime para \texttt{pptxgenjs}
  \item \textbf{LibreOffice} --- conversión pptx $\rightarrow$ pdf para QA visual
\end{itemize}

\subsection{Plataformas de despliegue y distribución}
\begin{itemize}
  \item \textbf{GitHub} --- \url{https://github.com/caledelta/UrbanFlow_CDMX}
  \item \textbf{Google Colab} --- ejecución del notebook reproducible
  \item \textbf{Streamlit} --- demo local con \texttt{streamlit run app/streamlit\_app.py}
\end{itemize}

\subsection{Recursos académicos y bibliográficos}
\begin{itemize}
  \item Norris, J. R. (1997). \textit{Markov Chains}. Cambridge University Press
  \item Rubinstein \& Kroese (2016). \textit{Simulation and the Monte Carlo Method} (3rd ed.). Wiley
  \item Van Lint \& Hoogendoorn (2010). A Robust and Efficient Method for Fusing Heterogeneous Data from Traffic Sensors. \textit{Computer-Aided Civil and Infrastructure Engineering}, 25(8)
  \item Harris et al. (2020). Array programming with NumPy. \textit{Nature}, 585
  \item Sinnott, R. W. (1984). Virtues of the Haversine. \textit{Sky and Telescope}, 68(2)
\end{itemize}

\subsection{Recursos de mercado y contexto}
\begin{itemize}
  \item Mordor Intelligence (2025). \textit{Mexico Last Mile Delivery Market Report 2025--2030}
  \item The Logistics World (2024). Logística de última milla: innovaciones para entregas rápidas en 2025
  \item Asociación Mexicana de Venta Online --- Estudio sobre Venta Online en México 2025
  \item T21 (2026). Última milla, el eslabón más crítico de la logística en México hacia 2030
\end{itemize}

% ============================================================
\section{Resumen ejecutivo del trabajo realizado}

\subsection{Métricas del proyecto}
\begin{itemize}
  \item Líneas de código: \texttt{[POR CONFIRMAR]} en \texttt{src/} y \texttt{app/}
  \item Tests automatizados: \textbf{443 tests pasando al 100\,\%}
  \item Cobertura de tests: \texttt{[POR CONFIRMAR]}\,\%
  \item Fuentes de datos integradas: \textbf{5}
  \item Perturbaciones en el catálogo: \textbf{11}
  \item Rutas en el dataset de validación: \textbf{30}
  \item Páginas del artículo técnico: \textbf{23}
  \item Slides de la presentación: \textbf{16}
  \item Celdas del notebook Colab: \textbf{38}
\end{itemize}

\subsection{Resultados empíricos}
\begin{itemize}
  \item MAPE del $P_{50}$ (VialAI): \textbf{11.9\,\%}
  \item MAPE del ETA (TomTom): 19.0\,\%
  \item Ventaja de VialAI sobre TomTom: \textbf{7.1 puntos porcentuales}
  \item Cobertura empírica de la banda: \textbf{93.3\,\%}
  \item MAE en minutos: \textbf{4.8 min}
  \item Latencia del motor: \textbf{$<$500\,ms} por consulta
\end{itemize}

\subsection{Entregables finales}
\begin{enumerate}
  \item \texttt{UrbanFlow\_CDMX\_Articulo\_v2.pdf} --- artículo técnico (23 pp)
  \item \texttt{UrbanFlow\_CDMX\_Colab.ipynb} --- notebook reproducible
  \item \texttt{UrbanFlow\_CDMX\_Presentacion.pptx} --- presentación ejecutiva (16 slides)
  \item \texttt{VialAI\_Pitch\_B2B.pdf} --- pitch comercial para última milla (4 pp)
  \item \texttt{guion\_video.md} --- guion del video de 5 minutos
  \item \texttt{checklist\_demo.md} --- checklist para demo en vivo
  \item \texttt{Inventario\_Proyecto.pdf} --- este documento
  \item Código fuente completo en \url{https://github.com/caledelta/UrbanFlow_CDMX}
\end{enumerate}

\vfill
\begin{center}
{\color{navy}\rule{0.6\textwidth}{0.5pt}}\\[0.5em]
\textit{Documento preparado por} \textbf{Carlos Armando López Encino}\\
Diplomado en Ciencia de Datos \(\cdot\) FES Acatlán \(\cdot\) UNAM\\
Abril 2026
\end{center}

\end{document}
